\subsection{Directed Graphs and Interventions}
% DAGs, Markov factorization, d-separation, Markov equivalence, estimation, and intervention calculus.

Directed graphs encode conditional-independence structure with arrows. For causal use, arrows are not merely predictive: they specify how the data-generating process changes under intervention.

\begin{definition} \label{def:directed-graph-dag-language} \hlt{Directed Graphs and DAGs}
\begin{enumerate}[label=(\roman*)]
\setlength{\itemsep}{0pt}
\item A \hlt{directed graph} is $\mathcal{G}=(V,E)$ with vertices $V$ and directed edges $E\subseteq V\times V$.
\item If $(X,Y)\in E$, write $X\to Y$; $X$ is a \hlt{parent} of $Y$, $Y$ is a \hlt{child} of $X$, and $\pi(X)$ denotes the parents of $X$.
\item A \hlt{directed path} follows arrows; an \hlt{undirected path} ignores arrow directions.
\item If $X$ has a directed path to $Y$, then $X$ is an \hlt{ancestor} of $Y$ and $Y$ is a \hlt{descendant} of $X$.
\item A \hlt{cycle} is a directed path returning to its start; a graph with no cycles is a \hlt{directed acyclic graph}, or \hlt{DAG}.
\end{enumerate}
\end{definition}

\begin{definition} \label{def:collider-and-noncollider} \hlt{Colliders}
\begin{enumerate}[label=(\roman*)]
\setlength{\itemsep}{0pt}
\item On a three-node path $X-Y-Z$, the middle node $Y$ is a \hlt{collider} iff the local orientation is $X\to Y\leftarrow Z$.
\item The local orientations $X\to Y\to Z$, $X\leftarrow Y\leftarrow Z$, and $X\leftarrow Y\to Z$ are \hlt{non-colliders} at $Y$.
\end{enumerate}
\end{definition}

\begin{definition} \label{def:dag-markov-factorization} \hlt{Markov Factorization for a DAG}
\begin{enumerate}[label=(\roman*)]
\setlength{\itemsep}{0pt}
\item Let $\mathcal{G}$ be a DAG on $V=(X_1,\ldots,X_k)$, and let $\pi_i$ denote the values of the parents of $X_i$.
\item A mass or density $p$ is \hlt{Markov to $\mathcal{G}$} iff $p(x_1,\ldots,x_k)=\prod_{i=1}^{k}p(x_i\mid\pi_i)$.
\item The set of distributions Markov to $\mathcal{G}$ is $\mathcal{M}(\mathcal{G})$.
\end{enumerate}
\end{definition}

\begin{example} \label{ex:dag-factorization-v-structure-tail} \hlt{A Four-Node Factorization}\\
For the DAG $X\to Z\leftarrow Y$ and $Z\to W$, the Markov factorization is $p(x,y,z,w)=p(x)p(y)p(z\mid x,y)p(w\mid z)$.\\
It implies $X\perp\!\!\!\perp Y$ and $W\perp\!\!\!\perp\{X,Y\}\mid Z$.
\end{example}

\begin{theorem} \label{thm:dag-markov-condition} \hlt{Markov Condition}\\
Let $\mathcal{G}$ be a DAG on $V=(X_1,\ldots,X_k)$ in topological order.\\
For each $i$, let $P_i\defeq\operatorname{pa}_{\mathcal{G}}(X_i)$, $D_i\defeq\operatorname{de}_{\mathcal{G}}(X_i)$, and $N_i\defeq V\setminus(\{X_i\}\cup D_i)$.\\
For any joint mass or density $p$, $p\in\mathcal{M}(\mathcal{G})$ iff $X_i\perp\!\!\!\perp N_i\setminus P_i\mid P_i$ for every $i$.
\begin{proof}
\begin{enumerate}[label=\roman*.]
\setlength{\itemsep}{0pt}
\item[$(\Rightarrow)$] Assume $p(v)=\prod_{j=1}^{k}p(x_j\mid\pi_j)$. Fix $i$ and set $A_i\defeq N_i\setminus P_i$.\\
Marginalising descendants of $X_i$ in reverse topological order removes their conditional factors, giving $p(x_i,a_i,\pi_i)=p(x_i\mid\pi_i)c(a_i,\pi_i)$. Hence $p(x_i\mid a_i,\pi_i)=p(x_i\mid\pi_i)$, so $X_i\perp\!\!\!\perp A_i\mid P_i$.
\item[$(\Leftarrow)$] Assume $X_i\perp\!\!\!\perp N_i\setminus P_i\mid P_i$ for every $i$.\\
Since the order is topological, $X_{<i}\subseteq N_i$, hence $p(x_i\mid x_{<i})=p(x_i\mid\pi_i)$.\\
The chain rule gives $p(v)=\prod_{i=1}^{k}p(x_i\mid x_{<i})=\prod_{i=1}^{k}p(x_i\mid\pi_i)$, so $p\in\mathcal{M}(\mathcal{G})$.
\end{enumerate}
\hlt{Structural Commentary}: The global factorization is equivalent to local parent-screening conditions.\\
\hlt{Key Insight}: In a topological order, predecessors are non-descendants.
\end{proof}
\end{theorem}

\begin{example} \label{ex:five-node-markov-condition} \hlt{Independences from a DAG}\\
For DAG $A\to B\to D\to E$, $A\to C\to D$, factorisation is $p(a,b,c,d,e)=p(a)p(b\mid a)p(c\mid a)p(d\mid b,c)p(e\mid d)$.\\
The Markov condition gives $B\perp\!\!\!\perp C\mid A$, $D\perp\!\!\!\perp A\mid\{B,C\}$, and $E\perp\!\!\!\perp\{A,B,C\}\mid D$.
\end{example}

\begin{figure}[H]
  \centering
  % Source: attached directed graph screenshots, Figures 20.5 and 20.6; redrawn as TikZ.
\begin{tikzpicture}[
  >=Stealth,
  line cap=round,
  line join=round,
  every node/.style={font=\small},
  v/.style={inner sep=1pt},
  cond/.style={
    draw=blue!60!black,
    fill=blue!6,
    rounded corners=1pt,
    inner xsep=2.5pt,
    inner ysep=1.5pt
  },
  edge/.style={->, black!78, line width=0.85pt}
]
  \begin{scope}
    \node[v] (x1) at (0,0) {$X$};
    \node[cond] (y1) at (0.9,0) {$Y$};
    \node[v] (z1) at (1.8,0) {$Z$};
    \draw[edge] (x1) -- (y1);
    \draw[edge] (y1) -- (z1);
  \end{scope}

  \begin{scope}[xshift=2.55cm]
    \node[v] (x2) at (0,0) {$X$};
    \node[cond] (y2) at (0.9,0) {$Y$};
    \node[v] (z2) at (1.8,0) {$Z$};
    \draw[edge] (y2) -- (x2);
    \draw[edge] (z2) -- (y2);
  \end{scope}

  \begin{scope}[xshift=5.10cm]
    \node[v] (x3) at (0,0) {$X$};
    \node[cond] (y3) at (0.9,0) {$Y$};
    \node[v] (z3) at (1.8,0) {$Z$};
    \draw[edge] (y3) -- (x3);
    \draw[edge] (y3) -- (z3);
  \end{scope}

  \begin{scope}[xshift=7.65cm]
    \node[v] (x4) at (0,0) {$X$};
    \node[cond] (y4) at (0.9,0) {$Y$};
    \node[v] (z4) at (1.8,0) {$Z$};
    \draw[edge] (x4) -- (y4);
    \draw[edge] (z4) -- (y4);
  \end{scope}

  \begin{scope}[xshift=10.20cm]
    \node[v] (x5) at (0,0) {$X$};
    \node[v] (y5) at (0.9,0) {$Y$};
    \node[v] (z5) at (1.8,0) {$Z$};
    \node[cond] (w5) at (0.9,-0.72) {$W$};
    \draw[edge] (x5) -- (y5);
    \draw[edge] (z5) -- (y5);
    \draw[edge] (y5) -- (w5);
  \end{scope}
\end{tikzpicture}

  \caption{Chain, reverse chain, fork, collider, descendant collider; blue = conditioned.}
  \label{fig:d-separation-motifs}
\end{figure}

\begin{definition} \label{def:d-separation} \hlt{D-Separation}
\begin{enumerate}[label=(\roman*)]
\setlength{\itemsep}{0pt}
\item Let $X\neq Y$ be vertices and let $W\subseteq V\setminus\{X,Y\}$.
\item An undirected path $\gamma:X\leadsto Y$ is \hlt{active given $W$} iff every collider on $\gamma$ has a descendant in $W$ and no non-collider on $\gamma$ lies in $W$.
\item Vertices $X,Y$ are \hlt{d-separated given $W$} iff no path $\gamma:X\leadsto Y$ is active given $W$.
\item If some path $\gamma:X\leadsto Y$ is active given $W$, then $X,Y$ are \hlt{d-connected given $W$}.
\end{enumerate}
\end{definition}

\begin{method} \label{meth:d-separation-rules} \hlt{Rules of D-Separation}\\
For a path segment with endpoints $X$ and $Z$ and middle node $Y$:
\begin{enumerate}[label=(\roman*)]
\setlength{\itemsep}{0pt}
\item If $Y$ is a non-collider, the path is active unless one conditions on $Y$.
\item If $Y$ is a collider, the path is blocked unless one conditions on $Y$ or on a descendant of $Y$.
\item Conditioning on a descendant of a collider opens the path just as conditioning on the collider does.
\end{enumerate}
\end{method}

\begin{example} \label{ex:d-separation-chain-fork-collider} \hlt{Three-Node Patterns}\\
For any non-collider triple $X-Y-Z$, endpoints $X,Z$ are d-connected given $\varnothing$ and d-separated given $Y$.\\
For the collider $X\to Y\leftarrow Z$, this reverses.
\end{example}

\begin{theorem} \label{thm:d-separation-conditional-independence-reading} \hlt{D-Separation and Conditional Independence}\\
Let $\mathcal{G}$ be a DAG on $V$, and let $p\in\mathcal{M}(\mathcal{G})$.
For disjoint $A,B,C\subseteq V$, if $A$ and $B$ are d-separated by $C$ in $\mathcal{G}$, then $A\perp\!\!\!\perp B\mid C$ under $p$.
If $\mathcal{I}(p)$ is the set of conditional independences under $p$ and $\mathcal{I}(\mathcal{G})$ is the set implied by d-separation, then faithfulness of $p$ to $\mathcal{G}$ gives $\mathcal{I}(p)=\mathcal{I}(\mathcal{G})$.
\begin{proof} Let $H$ be the ancestral subgraph induced by $\operatorname{an}(A\cup B\cup C)$, and let $H^m$ be its moral graph.\\
The d-separation criterion is equivalent to ordinary separation of $A$ and $B$ by $C$ in $H^m$.\\
Since $p\in\mathcal{M}(\mathcal{G})$, marginalizing variables outside $H$ and moralizing parent sets yields a factorization whose factors are supported on cliques of $H^m$.\\
If $C$ separates $A$ and $B$ in $H^m$, those clique factors partition as $\phi_A(a,c)\phi_B(b,c)$, so $p(a,b,c)=\phi_A(a,c)\phi_B(b,c)$ and $A\perp\!\!\!\perp B\mid C$.\\
Thus $\mathcal{I}(\mathcal{G})\subseteq\mathcal{I}(p)$, and faithfulness gives $\mathcal{I}(p)\subseteq\mathcal{I}(\mathcal{G})$, hence $\mathcal{I}(p)=\mathcal{I}(\mathcal{G})$.\\
\hlt{Structural Commentary}: Markovness gives graphical soundness; faithfulness adds completeness.\\
\hlt{Key Insight}: D-separation is the graph-theoretic language for conditional independence.
\end{proof}
\end{theorem}

\begin{example} \label{ex:collider-descendant-d-separation} \hlt{Opening a Collider Through Descendants}\\
In the DAG $X\to U\leftarrow V\to W\leftarrow Y$, $U\to S_1$, and $W\to S_2$, the vertices $X,Y$ are d-separated given $\varnothing$.\\
They are d-connected given $\{S_1,S_2\}$, since descendants of the colliders $U,W$ are conditioned on, and d-separated given $\{S_1,S_2,V\}$, since the non-collider $V$ is conditioned on.
\end{example}

\begin{definition} \label{def:markov-equivalence-dag} \hlt{Markov Equivalence}
\begin{enumerate}[label=(\roman*)]
\setlength{\itemsep}{0pt}
\item Let $\mathcal{I}(\mathcal{G})$ denote the conditional independence model implied by a DAG $\mathcal{G}$.
\item DAGs $\mathcal{G}_1,\mathcal{G}_2$ on the same vertex set are \hlt{Markov equivalent} iff $\mathcal{I}(\mathcal{G}_1)=\mathcal{I}(\mathcal{G}_2)$.
\item The \hlt{skeleton} of $\mathcal{G}$ is the undirected graph obtained by replacing each arrow with an undirected edge.
\end{enumerate}
\end{definition}

\begin{theorem} \label{thm:markov-equivalence-skeleton-colliders} \hlt{Markov Equivalence Criterion}\\
Let $\mathcal{G}_1,\mathcal{G}_2$ be DAGs on the same vertex set $V$.\\
Let $\operatorname{sk}(\mathcal{G})$ be the skeleton and let $\operatorname{ucol}(\mathcal{G})\defeq\{(a,b,c):a\to b\leftarrow c,\ a\not\sim c\}$ be the set of unshielded colliders.\\
Then $\mathcal{G}_1,\mathcal{G}_2$ are Markov equivalent iff $\operatorname{sk}(\mathcal{G}_1)=\operatorname{sk}(\mathcal{G}_2)$ and $\operatorname{ucol}(\mathcal{G}_1)=\operatorname{ucol}(\mathcal{G}_2)$.
\begin{proof}
\begin{enumerate}[label=\roman*.]
\setlength{\itemsep}{0pt}
\item[$(\Rightarrow)$] In any DAG, vertices $a,b$ are adjacent iff no $C\subseteq V\setminus\{a,b\}$ d-separates them.\\
If they are nonadjacent, the later vertex in a topological order is d-separated from the earlier vertex by its parents. Thus equality of d-separation models gives equal skeletons.\\
For unshielded $a-b-c$, $b$ is a collider iff every $C$ with $a\perp_{\mathcal{G}}c\mid C$ excludes $b$.\\
Hence equal d-separations give equal unshielded colliders.
\item[$(\Leftarrow)$] Suppose $\operatorname{sk}(\mathcal{G}_1)=\operatorname{sk}(\mathcal{G}_2)$ and $\operatorname{ucol}(\mathcal{G}_1)=\operatorname{ucol}(\mathcal{G}_2)$.\\
For any disjoint $A,B,C\subseteq V$, existence of a $C$-active path from $A$ to $B$ is determined by the skeleton and the unshielded colliders; shielded orientation differences can be bypassed along the chord.\\
Hence $A$ and $B$ are d-separated by $C$ in $\mathcal{G}_1$ iff they are d-separated by $C$ in $\mathcal{G}_2$, so $\mathcal{I}(\mathcal{G}_1)=\mathcal{I}(\mathcal{G}_2)$.
\end{enumerate}
\hlt{Structural Commentary}: Markov equivalence is graphical equality of d-separation models.\\
\hlt{Key Insight}: D-separation sees only adjacencies and unshielded colliders.
\end{proof}
\end{theorem}

\begin{example} \label{ex:three-node-markov-equivalence} \hlt{Three Markov-Equivalent Paths}\\
The non-colliders $X\to Y\to Z$, $X\leftarrow Y\leftarrow Z$, and $X\leftarrow Y\to Z$ share skeleton $X-Y-Z$ and are Markov equivalent; the collider $X\to Y\leftarrow Z$ is not.
\end{example}

\begin{method} \label{meth:discrete-dag-maximum-likelihood} \hlt{Maximum Likelihood for a Discrete DAG}
\begin{enumerate}[label=(\roman*)]
\setlength{\itemsep}{0pt}
\item Let $V=(X_1,\ldots,X_k)$ be discrete and let the DAG be fixed.
\item For node $X_i$, write $\widetilde X_i$ for its parent configuration.
\item Given IID data $V_1,\ldots,V_n$, set $\widehat p(v)=\prod_{i=1}^{k}\widehat p(x_i\mid \widetilde x_i)$, where\\
$\widehat p(x_i\mid \widetilde x_i)=\#\{r:X_{ri}=x_i,\widetilde X_{ri}=\widetilde x_i\}/\#\{r:\widetilde X_{ri}=\widetilde x_i\}$ when the denominator is positive.
\end{enumerate}
\end{method}

\begin{method} \label{meth:parametric-continuous-dag-likelihood} \hlt{Parametric Likelihood for a Continuous DAG}
\begin{enumerate}[label=(\roman*)]
\setlength{\itemsep}{0pt}
\item Let $V_r=(X_{r1},\ldots,X_{rm})$, $r=1,\ldots,n$, and let the DAG be fixed.
\item Specify nodewise parametric conditionals $p(x_j\mid \pi_j;\theta_j)$.
\item The likelihood factorizes as $\mathcal{L}(\theta)=\prod_{r=1}^{n}p(V_r;\theta)=\prod_{r=1}^{n}\prod_{j=1}^{m}p(X_{rj}\mid \pi_{rj};\theta_j)$.
\item Thus a fixed DAG reduces multivariate likelihood to nodewise conditional likelihoods.
\end{enumerate}
\end{method}

\begin{remark} \label{rem:dag-structure-learning-caution} \hlt{Structure Learning Caution}\\
Scores $S(\mathcal{G};D)$ based on likelihood or information criteria can rank DAGs, but reliable recovery of $\mathcal{G}$ is sample-intensive; here DAGs mainly encode $\mathcal{I}(\mathcal{G})$ and contextual causal assumptions.
\end{remark}

\begin{figure}[H]
  \centering
  % Source: attached directed graph screenshots, Figure 20.9 and intervention pseudo-code; redrawn as TikZ.
\begin{tikzpicture}[
  >=Stealth,
  line cap=round,
  line join=round,
  every node/.style={font=\small},
  v/.style={inner sep=1pt},
  cond/.style={
    draw=blue!60!black,
    fill=blue!6,
    rounded corners=1pt,
    inner xsep=3pt,
    inner ysep=1.5pt
  },
  edge/.style={->, black!78, line width=0.85pt},
  cut/.style={densely dotted, black!65, line width=0.85pt}
]
  \begin{scope}
    \node[v] (x1) at (0,0) {$X$};
    \node[cond] (y1) at (1.65,0) {$Y$};
    \node[v] (z1) at (3.30,0) {$Z$};
    \draw[edge] (x1) -- (y1);
    \draw[edge] (y1) -- (z1);
    \draw[edge] (x1) to[bend left=32] (z1);
  \end{scope}

  \begin{scope}[xshift=5.10cm]
    \node[v] (x2) at (0,0) {$X$};
    \node[cond] (y2) at (1.65,0) {$Y$};
    \node[v] (z2) at (3.30,0) {$Z$};
    \draw[edge] (y2) -- (z2);
    \draw[edge] (x2) to[bend left=32] (z2);
    \draw[cut] (x2) -- (y2);
  \end{scope}
\end{tikzpicture}

  \caption{Left observe $Y=y$: $\sum_x p(z\mid x,y)p(x\mid y)$; right set $Y=y$: $\sum_x p(z\mid x,y)p(x)$.}
  \label{fig:conditioning-versus-intervention}
\end{figure}

\begin{definition} \label{def:passive-and-active-conditioning} \hlt{Passive and Active Conditioning}
\begin{enumerate}[label=(\roman*)]
\setlength{\itemsep}{0pt}
\item Let $\mathcal{G}$ be $X\to Y\to Z$ and $X\to Z$, with $p(x,y,z)=p(x)p(y\mid x)p(z\mid x,y)$.
\item \hlt{Passive conditioning}: observing $Y=y$ gives $\mathbb{P}(Z=z\mid Y=y)=\sum_x p(z\mid x,y)p(x\mid y)$.
\item \hlt{Active conditioning}: setting $Y:=y$ removes the mechanism $p(y\mid x)$\\
and gives $\mathbb{P}(Z=z\mid Y:=y)=\sum_x p(z\mid x,y)p(x)$.
\end{enumerate}
\end{definition}

\begin{method} \label{meth:dag-intervention-truncated-factorization} \hlt{Intervention by Truncated Factorization}\\
For a DAG factorization $p(v)=\prod_i p(x_i\mid \pi_i)$, the intervention ``set $X_j=x$'' is represented by:
\begin{enumerate}[label=(\roman*)]
\setlength{\itemsep}{0pt}
\item deleting all arrows into $X_j$ in the graph;
\item deleting the factor $p(x_j\mid \pi_j)$ and fixing $X_j=x$ in the remaining factors.
\end{enumerate}
The modified graph and distribution represent the intervened system.
\end{method}

\begin{example} \label{ex:rain-wet-lawn-causal-direction} \hlt{Rain and Wet Lawn}\\
Let $R$ be rain and $W$ wet lawn. Both $R\to W$ and $R\leftarrow W$ can factor a dependent joint law, but causally $R\to W$ gives $p(r\mid W:=w)=p(r)$, whereas $R\leftarrow W$ would give the implausible $p(r\mid W:=w)=p(r\mid w)$.
\end{example}

\begin{figure}[H]
  \centering
  % Source: attached directed graph screenshots, Figure 20.10; redrawn as TikZ.
\begin{tikzpicture}[
  >=Stealth,
  line cap=round,
  line join=round,
  every node/.style={font=\small},
  v/.style={inner sep=1pt},
  obs/.style={
    draw=blue!60!black,
    fill=blue!6,
    rounded corners=1pt,
    inner xsep=3pt,
    inner ysep=1.5pt
  },
  latent/.style={
    draw=black!65,
    dashed,
    rounded corners=1pt,
    inner xsep=3pt,
    inner ysep=1.5pt
  },
  edge/.style={->, black!78, line width=0.85pt}
]

  \begin{scope}
    \node[v] (xr) at (0,0) {$X$};
    \node[v] (yr) at (2.05,0) {$Y$};
    \node[v] (zr) at (1.02,1.05) {$Z$};
    \draw[edge] (xr) -- (yr);
    \draw[edge] (zr) -- (yr);
  \end{scope}

  \begin{scope}[xshift=3.05cm]
    \node[v] (xo) at (0,0) {$X$};
    \node[v] (yo) at (2.05,0) {$Y$};
    \node[obs] (zo) at (1.02,1.05) {$Z$};
    \draw[edge] (xo) -- (yo);
    \draw[edge] (zo) -- (xo);
    \draw[edge] (zo) -- (yo);
  \end{scope}

  \begin{scope}[xshift=6.10cm]
    \node[v] (xu) at (0,0) {$X$};
    \node[v] (yu) at (2.05,0) {$Y$};
    \node[latent] (zu) at (1.02,1.05) {$Z$};
    \draw[edge] (xu) -- (yu);
    \draw[edge] (zu) -- (xu);
    \draw[edge] (zu) -- (yu);
  \end{scope}
\end{tikzpicture}

  \caption{Left to right: randomized, observed $Z$, unobserved $Z$.}
  \label{fig:confounding-dag-types}
\end{figure}

\begin{proposition} \label{prop:dag-confounding-adjustment-connection} \hlt{DAG View of Confounding}\\
Let $Z$ be a common cause of treatment $X$ and outcome $Y$, with observational factorization\\
$p(z,x,y)=p(z)p(x\mid z)p(y\mid x,z)$, and set $m(x,z)\defeq\mathbb{E}(Y\mid X=x,Z=z)$.
\begin{enumerate}[label=(\roman*)]
\setlength{\itemsep}{0pt}
\item Observational conditioning gives $\mathbb{E}(Y\mid X=x)=\int m(x,z)\,dF_{Z\mid X=x}(z)$.
\item Intervention gives $\mathbb{E}(Y\mid X:=x)=\int m(x,z)\,dF_Z(z)$.
\item Under randomization, $X\perp\!\!\!\perp Z$, hence $\mathbb{E}_r(Y\mid X=x)=\mathbb{E}(Y\mid X:=x)$,\\
where $\mathbb{E}_r$ is expectation under the randomized graph.
\item If $Z$ is unobserved, $\mathbb{E}(Y\mid X:=x)$ is generally not identifiable from $p(x,y)$ alone.
\end{enumerate}
\begin{proof}
\begin{enumerate}[label=\roman*.]
\setlength{\itemsep}{0pt}
\item[(i)] Under the observational law $P$, $P(dy,dz\mid X=x)=P(dy\mid X=x,Z=z)\,dF_{Z\mid X=x}(z)$,\\
hence $\mathbb{E}(Y\mid X=x)=\int m(x,z)\,dF_{Z\mid X=x}(z)$.
\item[(ii)] Under intervention $P^x$, the truncated law is $P^x(dy,dz)=P(dy\mid X=x,Z=z)\,dF_Z(z)$,\\
hence $\mathbb{E}(Y\mid X:=x)=\int m(x,z)\,dF_Z(z)$.
\item[(iii)] Under randomization, $P^r(dz,dx,dy)=dF_Z(z)\,dF^r_X(x)\,P(dy\mid X=x,Z=z)$,\\
so $F^r_{Z\mid X=x}=F_Z$ and $\mathbb{E}_r(Y\mid X=x)=\int m(x,z)\,dF_Z(z)$.
\item[(iv)] If $Z$ is unobserved, $p(x,y)=\int p(z)p(x\mid z)p(y\mid x,z)\,dz$ is only a marginal mixture.\\
Different latent mixtures can share $p(x,y)$ but yield different $\int m(x,z)\,dF_Z(z)$.
\end{enumerate}
\hlt{Structural Commentary}: Confounding is a change in the mixing law for $Z$ when conditioning on $X$.\\
\hlt{Key Insight}: Conditioning averages over $F_{Z\mid X=x}$; intervention averages over $F_Z$.
\end{proof}
\end{proposition}

\begin{remark} \label{rem:counterfactual-types-as-confounders} \hlt{Counterfactual Types as Confounders}\\
For binary treatment $X$, outcome $Y$, $C_x\in\{0,1\}$ potential outcome under $X:=x$, set $Z\defeq(C_0,C_1)\in\{0,1\}^2$.\\
If $Z$ affects both treatment assignment and outcome, then $Z$ is DAG confounder encoding latent response type.
\end{remark}

\begin{example} \label{ex:discrete-conditioning-versus-intervention-logistic} \hlt{Discrete Conditioning Versus Intervention}\\
Let $\ell(t)=e^t/(1+e^t)$, $X\sim\operatorname{Bernoulli}(1/2)$, $Y\mid X=x\sim\operatorname{Bernoulli}(\ell(4x-2))$, $Z\mid X=x$,\\
$Y=y\sim\operatorname{Bernoulli}(\ell(2(x+y)-2))$. Observing $Y=y$ reweights the two possible $X$ values by $\mathbb{P}(X=x\mid Y=y)$, whereas setting $Y:=y$ keeps their prior weights $1/2$.\\
For $y=1$, with $a=\ell(-2)$ and $b=\ell(2)$, $\mathbb{P}(Z=1\mid Y=1)=a/2+b^2$ but $\mathbb{P}(Z=1\mid Y:=1)=1/4+b/2$.
\end{example}

\begin{example} \label{ex:gaussian-conditioning-versus-intervention} \hlt{Gaussian Structural Equations}\\
Let $X\sim N(0,1)$, $Y=\alpha X+\varepsilon_Y$, and $Z=\beta Y+\gamma X+\varepsilon_Z$, with independent standard normal errors. Since $\mathbb{E}(X\mid Y=y)=\alpha y/(1+\alpha^2)$, ordinary conditioning gives $Z\mid Y=y\sim N((\beta+\gamma\alpha/(1+\alpha^2))y,1+\gamma^2/(1+\alpha^2))$.\\
Intervention $Y:=y$ leaves $X\sim N(0,1)$ and gives $Z\mid Y:=y\sim N(\beta y,1+\gamma^2)$.
\end{example}

\begin{proposition} \label{prop:gaussian-confounding-correlation-can-mislead} \hlt{Marginal Correlation Can Mislead}\\
Let $X,\varepsilon_Y,\varepsilon_Z$ be independent $N(0,1)$ variables and let $Y=\alpha X+\varepsilon_Y$ and $Z=\beta Y+\gamma X+\varepsilon_Z$.\\
Then $\operatorname{Cov}(Y,Z)=\beta(1+\alpha^2)+\alpha\gamma$.
\begin{proof}
Since $\operatorname{Var}(Y)=\operatorname{Var}(\alpha X+\varepsilon_Y)=1+\alpha^2$ and $\operatorname{Cov}(Y,X)=\alpha$,\\
linearity gives $\operatorname{Cov}(Y,Z)=\operatorname{Cov}(Y,\beta Y+\gamma X+\varepsilon_Z)=\beta(1+\alpha^2)+\alpha\gamma$.
Thus $\operatorname{Cov}(Y,Z)\neq0$ is possible when $\beta=0$, while $\operatorname{Cov}(Y,Z)=0$ with $\beta\neq0$ occurs when $\alpha\neq0$ and $\gamma=-\beta(1+\alpha^2)/\alpha$.\\
\hlt{Structural Commentary}: Marginal covariance mixes direct effect and common-cause paths.\\
\hlt{Key Insight}: Correlation is not the structural coefficient $\beta$ when $\gamma X$ is present.
\end{proof}
\end{proposition}

\begin{example} \label{ex:randomized-gaussian-corrects-correlation-reading} \hlt{Randomization Restores the Correlation Reading}\\
Let $X\sim N(0,1)$, let $Y\sim N(\alpha,1)$ independently of $X$, and let $Z=\beta Y+\gamma X+\varepsilon_Z$ with $\varepsilon_Z\sim N(0,1)$ independent.\\
Then $\operatorname{Cov}(Y,Z)=\beta$ and $Z\mid Y:=y\sim N(\beta y,1+\gamma^2)$.\\
Thus zero randomized correlation is equivalent to no interventional dependence on $y$, namely $\beta=0$.
\end{example}
